%% P29 --- Entropia krzyżowa i dywergencja Kullbacka--Leiblera
%%
%% Zalążek. Poniższy opis jest kontraktem tego programu: co ma uzasadnić i co
%% ma dać czytelnikowi. Napisanie programu oznacza usunięcie bloku
%% \programstub{} i zastąpienie go programem. Jeśli po przerobieniu materiału
%% opis okaże się błędny, popraw go i odnotuj sprzeczność w CLAUDE.md w sekcji
%% *Resolved questions*.

\program{Entropia krzyżowa i dywergencja Kullbacka--Leiblera}
\label{prog:P29}

\programstub{%
\textit{[Opis do przetłumaczenia --- patrz wydanie angielskie.]}

Argument: cross-entropy is the cost of coding one distribution with
another's code; KL is the excess; it is non-negative, zero only when the
distributions agree, and it is not symmetric. Payoff: the question the brief
asks for --- what the asymmetry costs you when you pick a loss. Forward KL
is mode-covering: it pays an unbounded price for putting no mass where the
target has some, so it spreads. Reverse KL is mode-seeking: it pays for
putting mass where the target has none, so it collapses onto one mode.
Distillation minimises one, a variational objective and a KL-regularised
policy the other, and the difference is visible in the output. Measured on a
bimodal target, not asserted. Also: \enquote{KL distance} is a misnomer and the
triangle inequality fails; Jensen--Shannon as the symmetric alternative and
what it costs; and the observation that minimising cross-entropy against a
fixed dataset is minimising forward KL to the empirical distribution, which
closes the loop with P25.

\medskip
\textbf{Planowana długość:} 55 ramek.
}
